%%%%%%%%%%%%%%%%%
% 郭承宇 — 个人简历
% 基于 MAltaCV 模板 (github.com/Manciukic/MAltaCV)
% 编译方式: XeLaTeX → PDF
% 需要: maltacv.cls (同目录), 微软雅黑字体 (Windows 自带)
%%%%%%%%%%%%%%%%%

\documentclass[10pt,a4paper,ragged2e]{maltacv}

%% ===== 字体设置（经典正式：宋体 + Times New Roman）=====
\usepackage{fontspec}
\usepackage{xeCJK}

% 中文正文：宋体（衬线，正式）；加粗用黑体（宋体无粗体字形）
% 不设斜体（宋体无斜体，强设会变伪斜体歪扭）
\setCJKmainfont{SimSun}[BoldFont=SimHei]
% 显式声明：宋体斜体回退为正体（宋体无斜体字形，消除 Font Warning）
\setCJKmainfont{SimSun}[BoldFont=SimHei,ItalicFont=SimSun]
\setCJKsansfont{SimHei}
\setCJKmonofont{FangSong}        % 中文代码注释用仿宋

% 英文/数字：Times New Roman（与宋体搭配的经典组合，字形协调）
\setmainfont{Times New Roman}
\setsansfont{Times New Roman}
% 代码：等宽字体 Consolas（保证 INSERT...ON DUPLICATE 等代码对齐）
\setmonofont{Consolas}[Scale=0.85]

%% 抑制 biblatex Biber 警告（简历未使用参考文献）
\makeatletter
\def\blx@bcf@error#1{}
\makeatother

\usepackage{multicol}
\usepackage{setspace}
\setstretch{1.06}

%% 消除分栏多余的垂直间距
\setlength\multicolsep{0pt}

%% 调整列表左边距（覆盖模板默认值，适配中文）
\setlist[itemize]{leftmargin=2em,labelsep=0.5em,nosep,itemsep=1pt,label=\itemmarker}

%% ===== 色彩方案 =====
\setcolorscheme{raisinblack_flame}

%% ===== 姓名 =====
\name{郭承宇}
\firstname{郭承宇}
\familyname{}

\begin{document}

%% ===== 副标题 =====
\tagline{后端开发工程师}

%% ===== 个人信息栏 =====
\personalinfo{%
  \email{gcy\_4708@qq.com}
  \phone{15047894708}
  \github{4708gcy}
}

%% ===== 个人概述（留空：MAltaCV 的 \makecvheader 内部会读取 \@bio，
%%      不赋值会报 Undefined control sequence。赋空值即可：既不显示文字，又不破坏模板）=====
\bio{}

\makecvheader

%% ============================================
%% 技能
%% ============================================
\cvsection{专业技能}
\medskip

\begin{center}
  \begin{multicols}{4}
    \cvlistitem{C/C++17}{}
    \cvlistitem{Python}{}
    \cvlistitem{MySQL}{}
    \cvlistitem{FastAPI}{}
    \cvlistitem{Elasticsearch}{}
    \cvlistitem{Git \& CMake}{}
    \cvlistitem{数据结构与算法}{}
  \end{multicols}
\end{center}

%% ============================================
%% 教育背景
%% ============================================
\cvsection{教育背景}
\medskip

\cvuniversity
{软件工程（本科）}
{吉林大学（985 / 双一流）}
{2023.08 -- 2028.06（预计）}
{长春}

\begin{itemize}
  \item GPA：3.56 / 4.0 \quad | \quad 专业排名：40 / 350（前 11.4\%）
  \item 主修：数据库系统、操作系统、计算机网络、数据结构、算法设计与分析
\end{itemize}

%% ============================================
%% 项目经验
%% ============================================
\cvsection{项目经验}
\medskip

\cvexperience
{独立开发}
{HIS 医院信息系统}
{2026.04 -- 2026.06}
{课程设计项目}
{C++17, MySQL 8.0, FastAPI, LangChain, WinSock2 \quad \href{https://github.com/4708gcy/HIS}{github.com/4708gcy/HIS}}

\begin{itemize}
  \item \textbf{数据库架构设计：}
        独立设计 13 主表 + 7 子表的 MySQL 关系模型，含外键约束与 CHECK 约束；
        编写 Python 脚本完成 CSV 历史数据清洗与迁移

  \item \textbf{内存-数据库双向同步：}
        13 条双向链表作运行时数据结构，退出时每条链表独立事务包装
        \texttt{INSERT ... ON DUPLICATE KEY UPDATE} 与 MySQL 全量同步，异常自动回滚

  \item \textbf{架构重构与泛型抽象：}
        将 Admin 模块从单文件拆为 4 个职责子模块；
        抽象 \texttt{EntityRepository<T>} 泛型模板统一 3 条链表 O(1) ID 查找，
        \texttt{displayChainByFilter<T>} 消除遍历重复代码

  \item \textbf{安全机制：}
        手写 SHA-256 加盐哈希（10,000 次迭代，兼容旧版 1,000 次）；
        统一参数化封装 + 字符串转义防 SQL 注入；复方药物成分级分词匹配，
        解决"阿莫西林"误匹配"阿莫西林克拉维酸钾"假阳性

  \item \textbf{AI 分析服务（Python）：}
        基于 FastAPI + LangChain 搭建 9 个 REST 端点；传统统计
        （Holt-Winters 双指数预测、Z-score 异常检测）与 LLM 双策略路由；
        RAG 知识库（BGE 嵌入 + FAISS 向量检索）
\end{itemize}

\medskip

\cvexperience
{AI 辅助开发}
{PaperLens 论文与课件学习智能体}
{2026.06}
{个人项目}
{LangChain, LangGraph, FastAPI, Elasticsearch, BGE-m3 \quad \href{https://github.com/4708gcy/PaperLens}{github.com/4708gcy/PaperLens}}

\begin{itemize}
  \item \textbf{AI 辅助开发实践：}
        系统学习大模型工程并完成大量代码练习后，让 AI 读取已掌握的代码库生成工程方案，
        再由 AI 编程实现。核心价值在驾驭 AI 做复杂工程——通过提供大量真实上下文
        压缩模型自由发挥空间以减少幻觉、设技术边界、主导架构决策与产出验收

  \item \textbf{双通路 + 三级素材降级：}
        论文精读走 ES 多路召回 RAG（BM25 + BGE-m3 + RRF + Cross-Encoder 重排），
        学习助手走全文通读；超长文档按上下文预算自适应——通读优先，超限降级检索 / 截断，
        解决思考模型 98K 输入限制

  \item \textbf{LangGraph 多分支 Agent：}
        Send API 章节并行综述、多轮记忆注入历史窗口、思考模型分级（深度任务开 / 结构化关）。
        自建 60 条评测集对比 4 策略，LLM-as-Judge 验证 RAG 提升答案质量 +68.4\%
\end{itemize}

%% ============================================
%% 荣誉奖项
%% ============================================
\cvsection{荣誉奖项}
\medskip

\begin{multicols}{2}
  \cvlistitem{校级二等奖学金}{吉林大学}

  \vfill\null
  \columnbreak

  \cvlistitem{院优秀学生}{吉林大学}
\end{multicols}

\end{document}
